% Consolidated guide-selftest10 from 283f4ad03ab2102bd6bfdf75a06052da75ec7dbf; source fragments preserved.

% === BEGIN guide-selftest10-choices.tex blob ff1fa4508353b5e0238a6575ef1ad095c336dce8 ===
% 学习指导第5版自测题10（光的衍射与偏振）10.1.1—10.1.10；
% 原书印刷206—207，PDF物理214—215。每题按源选项和原图检查。
\subsection{自测题10：单缝、光栅与偏振选择题}

\sourcetest{10.1.1}{206}{214}
\problemtext 单缝光垂直入射，某$\theta=30^\circ$方向
按半波带法可以把波面分成3个半波带，
求单缝宽度$a$。
\workings
单缝两边光程差$\delta=a\sin\theta$。
每个半波带对应$\lambda/2$，故
$a\sin30^\circ=3\lambda/2$，即
\[
\boxed{a=3\lambda}.
\]
选\textbf{D}。
\crosscheck 3个半波带不是3个全波长，不能漏掉$1/2$。

\sourcetest{10.1.2}{206}{214}
\problemtext 夫琅禾费单缝衍射，狭缝宽增加为原来的$3/2$，
光波波长减至原来的$3/4$，问中央明纹近轴宽度比例。
\workings
中央明纹角宽或固定焦距的焦面线宽均正比$\lambda/a$，
\[
\boxed{\frac{W'}W=
\frac{(3/4)\lambda}{(3/2)a}\bigg/
\frac{\lambda}{a}=\frac12}.
\]
选\textbf{D}。
\crosscheck 中央明纹变窄，不是随缝宽增加而变宽。

\sourcetest{10.1.3}{206}{214}
\problemtext 原书图10.1.2：单缝宽$a$略增大，
并将整个单缝沿$+y$方向微小平移；
透镜和接收屏不动，问中央衍射纹宽度及位置变化。
\workings
单缝变宽使中央衍射明纹角宽
$\Delta\theta\simeq2\lambda/a$减小。
理想平面波与透镜后焦面构成傅里叶光学布置，
平移孔径仅使衍射场的复振幅多一个线性相位因子，
其强度模不变，中央亮纹不因单缝整体平移而移动。
所以\textbf{变窄，不移动}，选\textbf{C}。
\crosscheck 必须区分源题的“把缝宽稍微\emph{变宽}”
与“缝沿$+y$小位移”；不能把两种变化混作宽度减小。

\sourcetest{10.1.4}{206}{214}
\problemtext 已有光栅在屏幕上只能显示零级和正负一级主极大，
若想出现更高级主极大，应换大还是小的光栅常数？
\workings
光栅条件$d\sin\theta=m\lambda$，
传播方向满足$|m|\le d/\lambda$。
要增大允许最高级数$m_{\max}$，
需要更大的光栅常量$d$，
故选\textbf{B}。
\crosscheck 更大的刻线密度意味着更小的$d$，
反而减少可传播级次；“高分辨刻线多”与“高阶可见”是不同优化目标。

\sourcetest{10.1.5}{206}{214}
\problemtext 多缝光栅的所有偶数级主极大均遇单缝零点而缺失。
透光缝宽$a$与相邻不透明缝间隔$b$有何关系？
\workings
主极大$d\sin\theta=m\lambda$，
单缝暗纹$a\sin\theta=k\lambda$，
缺级当$m=(d/a)k$。
所有偶级缺失而奇级保留对应$d/a=2$，
即$d=a+b=2a$，所以
\[
\boxed{a=b}.
\]
选\textbf{A}。
\crosscheck “所有偶级”指非零偶数级；
零级中心主极大不在单缝零点。

\sourcetest{10.1.6}{207}{215}
\problemtext 强度$I_0$自然光依次通过两片理想偏振片，
透振轴夹角30°，求透射强度。
\workings
通过第一片：$I_1=I_0/2$。
再按马吕斯定律：
\[
\boxed{I=I_1\cos^230^\circ
=\frac{I_0}{2}\cdot\frac34=\frac{3I_0}{8}}.
\]
选\textbf{E}。
\crosscheck 首片透射因子$1/2$必不可少。

\sourcetest{10.1.7}{207}{215}
\problemtext 三片偏振片$P_1,P_2,P_3$，$P_1$与$P_3$互相正交。
自然光照射后出射光强为$I_0/8$。
只转动中间片$P_2$使强度归零，问最小转角。
\workings
设中片轴与$P_1$轴夹角$\theta$，$P_3$轴与$P_1$夹90°，
\[
 I=\frac{I_0}{2}\cos^2\theta\sin^2\theta
 =\frac{I_0}{8}\sin^22\theta.
\]
已知$I=I_0/8$，故$\sin^22\theta=1$，
中片轴在45°模90°位置。
要归零必须使其与某一片轴平行，最小转动
\[
\boxed{\Delta\theta_{\min}=45^\circ},
\]
选\textbf{B}。
\crosscheck 纯自然光通过首片后才能应用第二、第三片的马吕斯定律。

\sourcetest{10.1.8}{207}{215}
\problemtext 入射束混合自然光$I_n$与线偏振光$I_p$，
用理想旋转分析偏振片测得最大透射强度为最小的5倍。
求$I_n/I_p$。
\workings
自然光无论分析器转角，总有$I_n/2$通过；
线偏振分量平行分析器时全部通过、垂直时不通过。
于是
\[
\frac{I_{\max}}{I_{\min}}
=\frac{I_n/2+I_p}{I_n/2}=5
\quad\Longrightarrow\quad I_p=2I_n.
\]
故
\[
\boxed{\frac{I_n}{I_p}=\frac12}.
\]
选\textbf{A}。
\crosscheck 不应把自然光本身的半强度因子错施加到已完全线偏振的分量。

\sourcetest{10.1.9}{207}{215}
\problemtext 自然光以60°从空气入射透明介质，
测得反射光完全线偏振，求折射角以及折射光偏振性质。
\workings
布儒斯特条件要求
$i_B+r_B=90^\circ$，
所以
\[
\boxed{r_B=90^\circ-60^\circ=30^\circ}.
\]
入射自然光中平行入射面的$p$分量此时不反射，
但会折射入介质；$s$分量亦部分透射。
因此折射光一般仍为\textbf{部分偏振光}。
选\textbf{B}。
\crosscheck 完全偏振的是\emph{反射}光，而非折射光。

\sourcetest{10.1.10}{207}{215}
\problemtext 自然光以布儒斯特角从空气入射玻璃，
判断反射光电矢量的振动方向。
\workings
菲涅耳反射系数$p$分量在布儒斯特角为零，
$s$分量保留，故反射光\textbf{完全线偏振，电场垂直入射面}。
选\textbf{C}。
\crosscheck 题目问“振动方向”是指光电场方向，
不是光线传播方向；不能把两者混淆。
% === END guide-selftest10-choices.tex ===

% === BEGIN guide-selftest10-blanks.tex blob d6f12f380ab99d2d6aa59616383717c79b91d9c3 ===
% 《大学物理学学习指导》自测题10填空10.2.1—10.2.8；
% 原书印刷207—208，PDF物理215—216页。
\subsection{自测题10：衍射与偏振填空}

\sourcetest{10.2.1}{207}{215}
\problemtext 惠更斯原理引入什么概念？
菲涅耳用什么思想补充惠更斯原理？
\workings
惠更斯把波前上的每一点看作产生\textbf{子波}的源。
菲涅耳再引入\textbf{子波相干叠加／干涉}的思想，
使波前次级波的叠加能解释衍射明暗。
\crosscheck 仅使用惠更斯几何包络而没有相干相位，
不能定量计算干涉衍射暗纹。

\sourcetest{10.2.2}{207}{215}
\problemtext 夫琅禾费单缝衍射，屏上P为第二级暗纹。
问对应波面可划成几个半波带；
若单缝宽变成一半，P将成为哪一级明/暗纹？
\workings
原来第二级暗纹满足
$a\sin\theta_P=2\lambda$，
所以全缝光程差含4个半波带。
缝宽改为$a/2$时同方向
\[
\frac a2\sin\theta_P=\lambda,
\]
对应第一级暗纹。故填：
\[
\boxed{N_{\frac12\lambda}=4},\qquad
\boxed{m_{\rm dark}=1}.
\]
\crosscheck 源题固定屏幕P与透镜光学几何，
半缝后相位差减半，并非把屏上坐标减半。

\sourcetest{10.2.3}{208}{216}
\problemtext 可见光400—760nm垂直入射光栅，
哪一级可见光谱可以在不与紧邻高级光谱重叠的条件下
保持完整？
\workings
第$m$级的红端对应$m(760)$；
第$m+1$级紫端对应$(m+1)(400)$。
若完整第$m$级无重叠，要求
\[
m(760)<(m+1)(400)
\ \Longleftrightarrow\ 360m<400.
\]
正整数$m$仅有$m=1$满足，故
因此填\textbf{第一级光谱}。
\crosscheck 该判据针对相邻正级谱线在角位置的重叠，
还要求器件有足够的传播角范围才能观察其完整谱段。

\sourcetest{10.2.4}{208}{216}
\problemtext 多缝光栅常量$d=3\,\mu$m，透光缝宽$a=1\,\mu$m；
单缝衍射中央明纹内有多少条光栅主极大谱线？
\workings
光栅主极大$d\sin\theta=m\lambda$，
中央单缝包络界限$|a\sin\theta|<\lambda$。
故$|m|<d/a=3$，仅
$m=-2,-1,0,1,2$。
\[
\boxed{N=5}.
\]
\crosscheck $m=\pm3$正落单缝第一暗纹，不计入亮纹。

\sourcetest{10.2.5}{208}{216}
\problemtext 利用理想偏振片把一束入射线偏振光的振动方向旋转90°，
至少需要多少片偏振片？此时最大输出光强占输入的多少？
\workings
一片最终透振轴垂直入射振动方向的偏振片不能透光。
至少需一个中间轴和一个最终轴，即\textbf{两片}。
取中间角$\theta$，末片轴为90°，
\[
\frac{I_{\rm out}}{I_{\rm in}}=
\cos^2\theta\sin^2\theta
=\frac14\sin^22\theta\le\frac14.
\]
最大时$\theta=45^\circ$，
\[
\boxed{N_{\min}=2},\qquad\boxed{I_{\max}/I_{\rm in}=1/4}.
\]
\crosscheck 题目入射已经是线偏振光，
首片不需要额外乘自然光的$1/2$。

\sourcetest{10.2.6}{208}{216}
\problemtext 太阳光从空气射到平静水面$n=1.33$，
若反射光完全线偏振，求太阳光的仰角及反射光电矢量方向。
\workings
布儒斯特入射角对法线测量：
$i_B=\arctan n$。
太阳仰角是对水平水面测量，故
\[
\boxed{\alpha=90^\circ-\arctan(1.33)
\simeq36.94^\circ}.
\]
反射光仅有$s$偏振，
电场\textbf{垂直于入射面}。
源图中的入射面位于纸面内，故E矢量
\textbf{垂直纸面}。
\crosscheck 入射角与太阳仰角互余，不可将53.06°直接填作仰角。

\sourcetest{10.2.7}{208}{216}
\problemtext 图10.2.2展示五种光入射不同折射率界面的情况，
其中四幅为线偏振入射、一幅为自然光；
需在原图补出反射光与折射光，标注电矢量方向。
\workings
设入射面为图示纸面，
用$s$表示电场垂直入射面，
$p$表示电场平行入射面。
各向同性非磁界面的菲涅耳方程不会把两分量相互转换。
\begin{center}
\begin{tabular}{lll}
\hline
入射类型 & 反射光 & 透射光\\ \hline
$s$线偏振 & $s$线偏振 & $s$线偏振\\
$p$线偏振，$i\ne i_B$ & $p$线偏振 & $p$线偏振\\
$p$线偏振，$i=i_B$ & 无反射光 & $p$线偏振\\
自然光，$i\ne i_B$ & 部分偏振 & 部分偏振\\
自然光，$i=i_B$ & $s$线偏振 & 部分偏振\\
\hline
\end{tabular}
\end{center}
其中
$\boxed{i_B=\arctan(n_2/n_1)}$。
\crosscheck \textbf{原图逐线矢量重绘待二校：}
此处已给出决定源图A—E各图的物理分类准则，
但未将五个原版权子图重新描绘成可编辑原创TikZ示意，
因此本题的\emph{图形填绘部分未完成最终图解质检}。
不能把“原则正确”伪称每条箭头都已逐图校验。

\sourcetest{10.2.8}{208}{216}
\problemtext 在光学各向异性晶体内，
沿某特殊方向正常光与非常光有何相位传播特征？
该方向叫什么？只存在一个这种方向的晶体叫什么？
\workings
沿\textbf{光轴}方向，常光与非常光
的相位折射率相同、传播方向不再发生双折射分离。
只有一个光轴方向的晶体称为
因此填\textbf{单轴晶体}。
\crosscheck 光轴是晶体内\emph{特殊传播方向}，
不是指穿过晶体几何中心的一条固定直线；
非晶各向同性介质没有相同的双折射结构。
% === END guide-selftest10-blanks.tex ===

% === BEGIN guide-selftest10-solutions.tex blob 3c2eeb8cc54160118855a2a07658e7961bcdc09e ===
% 学习指导第5版自测题10计算10.3.1—10.3.4，
% 原书印刷208—209／源PDF216—217页。
\subsection{自测题10：衍射与偏振计算题}

\sourcetest{10.3.1}{208}{216}
\problemtext (1) 单缝宽$a=1.0\times10^{-2}$cm、
透镜焦距$f=50$cm，波长400nm、760nm的垂直入射光
所形成第一级\emph{侧明纹}中心的间距；
(2) 换成光栅常量$d=1.0\times10^{-3}$cm的光栅，
求两波长\emph{第一级主极大}中心间距。
\workings
首先$a=10^{-4}$m，$f=0.50$m。
单缝衍射\emph{第一级侧亮纹}在常用半波带近似下满足
$a\sin\theta\simeq(3/2)\lambda$，
故两波长同侧第一级亮纹中心的间距
\[
\boxed{\Delta x_{\rm slit,approx}
\simeq\frac{3f(\lambda_2-\lambda_1)}{2a}
=\frac{3(0.50)(360\times10^{-9})}{2(10^{-4})}
=2.70\,\mathrm{mm}}.
\]
若使用严格夫琅禾费第一侧极大根
$\tan\beta_1=\beta_1$，
$\beta_1\simeq4.493409$，
其光程系数$\beta_1/\pi\simeq1.43030$，
则
\[
\boxed{\Delta x_{\rm slit,exact}
\simeq\frac{\beta_1}{\pi}
\frac{f(\lambda_2-\lambda_1)}a
\simeq2.57\,\mathrm{mm}}.
\]
(2) 光栅常量$d=10^{-5}$m，
第一级主极大满足$d\sin\theta=\lambda$，故
\[
\boxed{\Delta x_{\rm grating}\simeq
\frac{f(\lambda_2-\lambda_1)}d
=\frac{0.50(360\times10^{-9})}{10^{-5}}
=18.0\,\mathrm{mm}}.
\]
\crosscheck \textbf{不能把单缝第一级侧明纹误当作第一暗纹。}
单缝侧亮纹的$3/2$是常用半波带近似，
严格峰值需解超越方程；光栅主极大方程本身不需要这一近似。

\sourcetest{10.3.2}{208}{216}
\problemtext 波长600nm光垂直入射多缝光栅，
第2级主极大位于衍射角30°，第3级缺级。
求(1)光栅常量$d=a+b$；
(2)可能的最小透光缝宽$a$；
(3)取该最小$a$后，$-90^\circ<\varphi<90^\circ$
内可见全部主极大的级次。
\workings
光栅方程$d\sin\varphi_m=m\lambda$，
\[
\boxed{d=\frac{2(600\,\mathrm{nm})}{\sin30^\circ}
=2400\,\mathrm{nm}=2.40\,\mu\mathrm m}.
\]
第3级缺级需$3\lambda/d=k\lambda/a$，
即$a=kd/3$。
实际透光缝应$0<a<d$，整数$k=1,2$，
故$a=0.80$或$1.60\,\mu$m，
\[
\boxed{a_{\min}=0.80\,\mu\mathrm m}.
\]
最小缝宽方案使$d/a=3$，所以非零3的倍数级缺失。
又$d/\lambda=4$，
严格角区间排除$|\varphi|=90^\circ$对应的$|m|=4$。
候选$m=0,\pm1,\pm2,\pm3$；
$\pm3$为缺级，最后
\[
\boxed{m=0,\pm1,\pm2}.
\]
合计\textbf{5条可见主亮纹}。
\crosscheck 光栅常量$d$与透光缝宽$a$不可混用；
若采用另一个允许缝宽$a=1.60\,\mu$m，
缺级结构须另算，不应仍给5条。

\sourcetest{10.3.3}{208--209}{216--217}
\problemtext 两片理想偏振片$P_1,P_2$叠置；
入射光由等强度自然光和线偏振光构成。
两次测量中，线偏振入射分量与$P_1$的角度
分别45°、30°；两片透振轴的夹角
第一次30°、第二次未知$\theta$。
第一次总透射强度为第二次的$3/4$。
求$\theta$及每次经过各片后的光强占\emph{原总入射光强}的比例。
\workings
设两种入射成分强度各为$J$，
故总入射强度$I_{\rm in}=2J$。
第一次过$P_1$：
\[
 I_{1,1}=\frac J2+J\cos^245^\circ=J,
 \qquad\boxed{I_{1,1}/I_{\rm in}=\frac12}.
\]
第一次再过$P_2$（夹角30°）：
\[
 I_{1,2}=J\cos^230^\circ=\frac34J,
 \qquad\boxed{I_{1,2}/I_{\rm in}=\frac38}.
\]
第二次过$P_1$：
\[
 I_{2,1}=\frac J2+J\cos^230^\circ
 =\frac54J,\qquad
 \boxed{I_{2,1}/I_{\rm in}=\frac58}.
\]
题给$I_{1,2}=(3/4)I_{2,2}$，
所以$I_{2,2}=J$。
此时$(5J/4)\cos^2\theta=J$，
\[
\boxed{\cos^2\theta=\frac45},\quad
\boxed{\theta=\arccos\sqrt{4/5}
\simeq26.565^\circ},\qquad
\boxed{I_{2,2}/I_{\rm in}=\frac12}.
\]
\crosscheck 线偏振与自然光各占一半，
但首片对二者透过因子不同；
第二次过首片的比例$5/8$不能漏掉自然光的$1/2$。

\sourcetest{10.3.4}{209}{217}
\problemtext 三层平行界面介质$I,II,III$，
折射率$n_1=1.33,n_2=1.50,n_3=1.00$；
自然光由介质I入射II。
若在I/II界面观察到完全线偏振反射光，
求入射角$i$；
第二个II/III界面的反射光是否也完全线偏振？
\workings
I/II的完全线偏振反射由布儒斯特条件给出
\[
\boxed{i=\arctan\frac{n_2}{n_1}
=\arctan\frac{1.50}{1.33}\simeq48.43^\circ}.
\]
此时在II中的折射角
$r=90^\circ-i\simeq41.57^\circ$，
由于两界面平行，这就是光在II/III界面的入射角。
第二个界面若也达布儒斯特角，则应有
\[
 i_{B,23}=\arctan\frac{n_3}{n_2}
 =\arctan(2/3)\simeq33.69^\circ.
\]
两角不等，故在II/III界面$p$反射系数一般不为零；
第二束反射光\textbf{并非完全线偏振，通常是部分偏振光}。
\crosscheck II到III的临界角
$i_c=\arcsin(n_3/n_2)\simeq41.81^\circ$，
源题传播方向$r\simeq41.57^\circ<i_c$，
并未达到全反射。
必须同时核对布儒斯特角和临界角，不能混淆。
% === END guide-selftest10-solutions.tex ===
